% TODO make hw stylesheet
\documentclass[11pt]{scrartcl}
\usepackage[a4paper, total={6in, 8in}]{geometry}
\usepackage[usenames,dvipsnames,svgnames]{xcolor}
\usepackage[shortlabels]{enumitem}
\usepackage[framemethod=TikZ]{mdframed}
\usepackage{amsmath,amssymb,amsthm}
\usepackage[colorlinks]{hyperref}
\usepackage{multicol}
\usepackage{tabularx}

\hypersetup{
	colorlinks=true,
	linkcolor=blue,
	filecolor=magenta,      
	urlcolor=magenta,
	pdftitle={MAT 362 HW}
}

\usepackage{microtype}
\usepackage{mathtools}
\usepackage[headsepline]{scrlayer-scrpage}
\usepackage{thmtools}
\usepackage[textsize=footnotesize]{todonotes}

\mdfdefinestyle{mdbluebox}{roundcorner=10pt,innerbottommargin=9pt,
	linecolor=blue,backgroundcolor=TealBlue!5,}
\declaretheoremstyle[headfont=\sffamily\bfseries\color{MidnightBlue},
mdframed={style=mdbluebox},]{thmbluebox}

\mdfdefinestyle{mdbloobox}{frametitlefont=\bfseries,innerbottommargin=8pt,
	nobreak=true,backgroundcolor=TealBlue!5,linecolor=blue,}
\declaretheoremstyle[headfont=\bfseries\color{MidnightBlue},
mdframed={style=mdbloobox},headpunct={\\[3pt]},postheadspace=0pt,]{thmbloobox}

\mdfdefinestyle{mdredbox}{frametitlefont=\bfseries,innerbottommargin=8pt,
	nobreak=true,backgroundcolor=Salmon!5,linecolor=RawSienna,}
\declaretheoremstyle[headfont=\bfseries\color{RawSienna},
mdframed={style=mdredbox},headpunct={\\[3pt]},postheadspace=0pt,]{thmredbox}

\mdfdefinestyle{mdgreenbox}{linecolor=ForestGreen,backgroundcolor=ForestGreen!5,
	linewidth=2pt,rightline=false,leftline=true,topline=false,bottomline=false,}
\declaretheoremstyle[headfont=\bfseries\sffamily\color{ForestGreen!70!black},
mdframed={style=mdgreenbox},headpunct={ --- },]{thmgreenbox}

\mdfdefinestyle{mdorangebox}{linecolor=RedOrange,backgroundcolor=RedOrange!5,
	linewidth=2pt,rightline=false,leftline=true,topline=false,bottomline=false,}
\declaretheoremstyle[headfont=\bfseries\sffamily\color{RedOrange!70!black},
mdframed={style=mdorangebox},headpunct={ --- },]{thmorangebox}

\mdfdefinestyle{mdblackbox}{linecolor=black,backgroundcolor=RedViolet!5!gray!5,
	linewidth=3pt,rightline=false,leftline=true,topline=false,bottomline=false,}
\declaretheoremstyle[mdframed={style=mdblackbox}]{thmblackbox}

\declaretheoremstyle[spaceabove=6pt,spacebelow=6pt]{basehead}

\declaretheorem[style=basehead,name=Problem]{problem}
\declaretheorem[style=basehead,name=Problem,numbered=no]{problem*}
\declaretheorem[style=thmbloobox,name=Setup]{setup} % for config geo
\declaretheorem[style=thmredbox,name=Example,sibling=problem]{example}
\declaretheorem[style=thmredbox,name=$\bigstar$ Problem,sibling=problem]{niceprob}
\declaretheorem[style=thmbluebox,name=Theorem,sibling=problem]{theorem}
\declaretheorem[style=thmbluebox,name=Corollary,sibling=theorem]{corollary}
\declaretheorem[style=thmbluebox,name=Proposition,sibling=theorem]{prop}
\declaretheorem[style=thmbluebox,name=Lemma,numberwithin=problem]{lemma}
\declaretheorem[style=thmbluebox,name=Theorem,numbered=no]{theorem*}
\declaretheorem[style=thmbluebox,name=Lemma,numbered=no]{lemma*}
\declaretheorem[style=thmgreenbox,name=Claim,numberwithin=problem]{claim}
\declaretheorem[style=thmgreenbox,name=Claim,numbered=no]{claim*}
\declaretheorem[style=thmblackbox,name=Remark,numberwithin=problem]{remark}
\declaretheorem[style=thmblackbox,name=Remark,numbered=no]{remark*}
\declaretheorem[style=thmgreenbox,name=Definition,sibling=theorem]{definition}
\declaretheorem[style=thmgreenbox,name=Definition,numbered=no]{definition*}
\declaretheorem[style=thmorangebox,name=Warning,sibling=theorem]{warning}
\declaretheorem[style=thmorangebox,name=Warning,numbered=no]{warning*}
\declaretheorem[style=thmorangebox,name=Note,numbered=no]{note}
\declaretheorem[style=thmblackbox,name=Exercise,sibling=problem]{exercise}
\declaretheorem[style=thmblackbox,name=Exercise,numbered=no]{exercise*}
\declaretheorem[style=thmblackbox,name=Question,sibling=problem]{question}
\declaretheorem[style=thmblackbox,name=Question,numbered=no]{question*}

\addtolength{\textheight}{3.14cm}
\providecommand{\ol}{\overline}
\providecommand{\ul}{\underline}
\providecommand{\eps}{\varepsilon}
\providecommand{\half}{\frac{1}{2}}
\providecommand{\dang}{\measuredangle} %% Directed angle
\providecommand{\CC}{\mathbb C}
\providecommand{\FF}{\mathbb F}
\providecommand{\NN}{\mathbb N}
\providecommand{\QQ}{\mathbb Q}
\providecommand{\RR}{\mathbb R}
\providecommand{\ZZ}{\mathbb Z}
\providecommand{\dg}{^\circ}
\providecommand{\ii}{\item}
\providecommand{\setdiff}{\smallsetminus}
\providecommand{\alert}[1]{{\sffamily\textbf{\textcolor{blue}{#1}}}}
\providecommand{\inv}{^{-1}}
\providecommand{\mb}[1]{\mathbf{#1}}
\providecommand{\dif}{\;d}

\reversemarginpar

\newenvironment{soln}{\begin{proof}[Solution]}{\end{proof}}
\newenvironment{walkthrough}{\noindent \textbf{Walkthrough.}}{\medskip}
\newlist{walk}{enumerate}{3}
\setlist[walk]{label=\bfseries (\alph*)}

\providecommand{\pow}{\operatorname{Pow}}
\providecommand{\vocab}[1]{{\textbf{\textcolor{ForestGreen}{#1}}}}
\providecommand{\lcm}{\operatorname{lcm}}
\providecommand{\abs}[1]{\left|#1\right|}

\usepackage{asymptote}
\begin{asydef}
	defaultpen(fontsize(10pt));
	unitsize(1cm);
	size(8cm); // set a reasonable default
	usepackage("amsmath");
	usepackage("amssymb");
	settings.tex="pdflatex";
	settings.outformat="pdf";
	// Replacement for olympiad+cse5 which is not standard
	import geometry;
	// recalibrate fill and filldraw for conics
	void filldraw(picture pic = currentpicture, conic g, pen fillpen=defaultpen, pen drawpen=defaultpen)
	{ filldraw(pic, (path) g, fillpen, drawpen); }
	void fill(picture pic = currentpicture, conic g, pen p=defaultpen)
	{ filldraw(pic, (path) g, p); }
	// some geometry
	pair foot(pair P, pair A, pair B) { return foot(triangle(A,B,P).VC); }
	pair orthocenter(pair A, pair B, pair C) { return orthocentercenter(A,B,C); }
	pair centroid(pair A, pair B, pair C) { return (A+B+C)/3; }
	// cse5 abbreviations
	path CP(pair P, pair A) { return circle(P, abs(A-P)); }
	path CR(pair P, real r) { return circle(P, r); }
	pair IP(path p, path q) { return intersectionpoints(p,q)[0]; }
	pair OP(path p, path q) { return intersectionpoints(p,q)[1]; }
	path Line(pair A, pair B, real a=0.6, real b=a) { return (a*(A-B)+A)--(b*(B-A)+B); }
	// cse5 more useful functions
	picture CC() {
		picture p=rotate(0)*currentpicture;
		currentpicture.erase();
		return p;
	}
	pair MP(Label s, pair A, pair B = plain.S, pen p = defaultpen) {
		Label L = s;
		L.s = "$"+s.s+"$";
		label(L, A, B, p);
		return A;
	}
	pair Drawing(Label s = "", pair A, pair B = plain.S, pen p = defaultpen) {
		dot(MP(s, A, B, p), p);
		return A;
	}
	path Drawing(path g, pen p = defaultpen, arrowbar ar = None) {
		draw(g, p, ar);
		return g;
	}
\end{asydef}

\usepackage{mathrsfs}

\ihead{\footnotesize MAT 362 HW01}
\ohead{\footnotesize Last compiled \today}

\title{\vspace{-2em}MAT 362 HW01}
\author{\large Aditya Pahuja}
\date{\large Due 02/06}
\begin{document}
\maketitle

\begin{problem*}1.2.2.
	Let $\alpha(t)$ be a parametrized curve which
	does not pass through the origin. If $\alpha(t_0)$ is the point
	of the trace of $\alpha$ which is closest to the origin
	and $\alpha'(t_0)\ne 0$, show that
	the position vector $\alpha(t_0)$ is orthogonal to $\alpha'(t_0)$.
\end{problem*}

\begin{soln}
	Consider the quantity $\|\alpha(t)\|^2 = \alpha(t)\cdot\alpha(t)$.
	The length of the vector $\alpha(t)$ is minimized at $t=t_0$,
	we have that the derivative of $\alpha(t)\cdot\alpha(t)$
	at $t=t_0$ is zero, namely
	\[2\alpha(t_0)\cdot\alpha'(t_0) = 0.\]
	The conditions in the problem guarantee that $\alpha(t_0)$
	and $\alpha'(t_0)$ are both not zero vectors, so they must be orthogonal
	in order for their dot product to be zero.
\end{soln}

\begin{problem*}1.2.4.
	Let $\alpha\colon I\to\RR$ be a parametrized curve
	and let $v\in\RR^3$ be a fixed vector. Assume that $\alpha'(t)$
	is orthogonal to $v$ for all $t\in I$ and that $\alpha(0)$
	is also orthogonal to $v$. Prove that $\alpha(t)$ is orthogonal to $v$
	for all $t\in I$.
\end{problem*}

\begin{soln}
	We can compute
	\[(\alpha(t)\cdot v)' = \alpha'(t)\cdot v = 0\]
	for all $t$, so $\alpha(t)\cdot v$ is constant on $I$.
	We already know that $\alpha(0)\cdot v = 0$,
	hence $\alpha(t)\cdot v = 0$ for all $t$.
\end{soln}

\begin{problem*}1.2.5.
	Let $\alpha\colon I\to\RR^3$ be a parametrized curve,
	with $\alpha'(t)\ne 0$ for all $t\in I$. Show that
	$|\alpha(t)|$ is a nonzero constant if and only if
	$\alpha(t)$ is orthogonal to $\alpha'(t)$ for all $t\in I$.
\end{problem*}

\begin{soln}
	Of course, $|\alpha(t)|$ is a nonzero constant if and only if its square
	$|\alpha(t)|^2 = \alpha(t)\cdot\alpha(t)$ is as well. We have
	\[\frac{d}{dt}|\alpha(t)|^2 = 2\alpha(t)\cdot\alpha'(t) = 0,\]
	so $\alpha(t)$ and $\alpha'(t)$ are orthogonal if and only if
	$\alpha(t)$ is nonzero and $t\mapsto |\alpha(t)|^2$ is a constant function
	$I\to\RR$, i.e. $|\alpha(t)|$ is a nonzero constant.
\end{soln}

\begin{problem*}1.3.5.
	Let $\alpha\colon(-1,+\infty)\to\RR^2$ be given by
	\[\alpha(t) = \left(\frac{3at}{1+t^3},\frac{3at^2}{1+t^3}\right).\]
	Prove that:
	\begin{enumerate}[(a)]
		\ii For $t = 0$, $\alpha$ is tangent to the $x$-axis.
		\ii As $t\to+\infty$, $\alpha(t)\to(0,0)$ and $\alpha'(t)\to(0,0)$.
		\ii Take the curve with the opposite orientation. Now, as $t\to-1$,
		the curve and its tangent approach the line $x+y+a=0$.
	\end{enumerate}
\end{problem*}

\begin{soln}
	(a) Of course, $\alpha(0) = (0,0)$, so the curve intersects
	the $x$-axis at $t=0$. To show that it's tangent to the $x$-axis,
	we can compute
	\[\alpha'(t) = \left(\frac{3a(1+t^3) - 3at(3t^2)}{(1+t^3)^2},
	\frac{6at(1+t^3) - 3at^2(3t^2)}{(1+t^3)^2}\right)
	= \left(\frac{3a-6at^3}{(1+t^3)^2},
	\frac{6at-3at^4}{(1+t^3)^2}\right),\]
	so that $\alpha'(0) = (3a,0)$ is a vector along the $x$-axis.
	
	(b) The coordinates of $\alpha$ and $\alpha'$ are all
	rational functions in $t$ such that in each case
	the degree of the polynomial in the numerator is smaller than
	that in the denominator. This is enough to imply that
	each of the rational functions goes to zero as $t$ grows
	arbitrarily large.
	
	(c) As $t\to-1$, the curve gets closer to $x+y+a=0$
	in the sense that $x+y$ becomes arbitrarily close to $-a$.
	In particular, the sum of coordinates of $\alpha(t)$ is
	\[\frac{3at(1+t)}{(1+t^3)} = \frac{3at}{1-t+t^2}\]
	which as $t\to-1$ approaches $\frac{-3a}{3} = -a$.
	Moreover, the tangent approaches $x+y=-a$ in that the derivative
	becomes arbitrarily close to a scalar multiple of $(1,-1)$,
	so that the slope of the tangent line approaches that of $x+y-a$.
	We can see this by noting that the sum of the coordinates of $\alpha'$,
	which is the derivative of $\frac{3at}{1-t+t^2}$, is equal to
	\[\frac{3a(1-t+t^2) - 3at(2t-1)}{(1-t+t^2)}\]
	which, by setting $t=-1$, we can compute is equal to zero
	in the limit as $t\to-1$.
\end{soln}

\begin{problem*}1.3.6.
	Let $\alpha(t) = (ae^{bt}\cos t,ae^{bt}\sin t)$ for $t\in\RR$
	with $a$ and $b$ constants, $a > 0$, and $b < 0$ be a parametrized curve.
	\begin{enumerate}[(a)]
		\ii Show that as $t\to +\infty$, $a(t)$ approaches the origin $0$,
		spiraling around it.
		\ii Show that $\alpha'(t)\to(0,0)$ as $t\to+\infty$ and that
		\[\lim_{t\to\infty} \int_{t_0}^t |\alpha'(u)|\dif u\]
		is finite.
	\end{enumerate}
\end{problem*}

\begin{soln}
	(a) Since $b < 0$,
	\[\lim_{t\to\infty} e^{bt} = 0,\]
	and $\cos t$ and $\sin t$ are bounded so $ae^{bt}\cos t$
	and $ae^{bt}\sin t$ both approach zero as $t$ goes to $+\infty$.
	
	(b) The tangent vector is
	\[\alpha'(t) =
	(abe^{bt}\cos(t) - ae^{bt}\sin t, abe^{bt}\sin t + ae^{bt}\cos t),\]
	and we have already shown that each of the addends in each coordinate
	will approach zero as $t$ grows large because the addends are all
	products of bounded expressions with $e^{bt}$;
	hence $\alpha'(t)$ approaches $(0,0)$ as $t\to+\infty$.
	Moreover, we can compute
	\[\int_{t_0}^t |\alpha'(u)|\dif u=
	\int_{t_0}^t \sqrt{a^2e^{2bu}(\cos^2u+\sin^2u)}\dif u
	= \int_{t_0}^t ae^{bu}\dif u = \frac{a}{b}(e^{bt} - e^{bt_0}),\]
	so, as $t\to+\infty$, the $e^{bt}$ term disappears, leaving behind
	the finite quantity
	\[\frac{-aeb^{t_0}}{b}\]
	(note that this is positive because $b < 0$).
\end{soln}

\begin{problem*}1.3.10.
	Let $\alpha\colon I\to\RR^3$ be a parametrized curve.
	Let $[a,b]\subseteq I$and set $\alpha(a) = p$, $\alpha(b) = q$.
	\begin{enumerate}[(a)]
		\ii Show that for any constant vector $v$, $|v| = 1$,
		\[(q-p)\cdot v = \int_a^b \alpha'(t)\cdot v\dif t
		\le \int_a^b |\alpha'(t)|\dif t.\]
		\ii Set $v = \frac{q-p}{|q-p|}$ and show that
		\[|\alpha(b) - \alpha(a)|\le \int_a^b |\alpha'(t)|\dif t;\]
		that is, the curve of shortest length from $a$ to $b$
		is the straight line joining these points.
	\end{enumerate}
\end{problem*}

\begin{soln}
	(a) The equality follows from the fundamental theorem of calculus,
	since, if $\alpha(t) = (x(t), y(t), z(t))$ and $v = (v_1, v_2, v_3)$, then
	\begin{align*}
		\int_a^b \alpha'(t)\cdot v\dif t
		&= \int_a^b x'(t)v_1 + y'(t)v_2 + z'(t)v_3\dif t\\
		&= (x(b)v_1 + y(b)v_2 + z(b)v_3) - (x(a)v_1 + y(a)v_2 + z(a)v_3)\\
		&= p\cdot v - q\cdot v = (p-q)\cdot v.
	\end{align*}
	Then, the inequality follows from
	\[\int_a^b \alpha'(t)\cdot v\dif t
	\le \int_a^b |\alpha'(t)\cdot v|\dif t
	\le \int_a^b |\alpha'(t)|\cdot|v|\dif t
	= \int_a^b |\alpha'(t)|\dif t,\]
	where $|\alpha'(t)\cdot v|\le|\alpha'(t)|\cdot|v|$ is an application
	of the Cauchy-Schwarz inequality and the last equality uses $|v| = 1$.
	
	(b) Taking $v$ to be the given expression shows that
	\[(q-p)\cdot v = \frac{(q-p)\cdot (q-p)}{|q-p|} =
	\frac{|q-p|^2}{|q-p|} = |q-p| = |\alpha(b)-\alpha(a)|\]
	can, as in the result of part (a),
	be bounded above by $\int_a^b |\alpha'(t)|\dif t$,
	which is what we wanted.
\end{soln}

\begin{problem*}1.4.5.
	Show that the equation of a plane passing through
	three noncollinear points $p_1$, $p_2$, and $p_3$ is given by
	\[(p - p_1)\wedge(p - p_2)\cdot(p - p_3) = 0,\]
	where $p$ is an arbitrary point in the plane.
\end{problem*}

\begin{soln}
	Recall that if three vectors $a_1$, $a_2$, $a_3$ are coplanar
	with the origin of $\RR^3$, then $\det(a_1,a_2,a_3) = 0$.
	If $p$ lies in the plane through $p_1$, $p_2$, $p_3$, then
	$p_1-p$, $p_2-p$, and $p_3-p$ are coplanar with $p-p$,
	since subtracting $p$ is just a translation by the vector $-p$.
	We can also negate all the vectors without affecting
	coplanarity, since reflecting through the origin doesn't affect
	coplanarity, so actually $p-p_1$, $p-p_2$, and $p-p_3$
	are coplanar with $0$. To finish, we have
	\[(p - p_1)\wedge(p - p_2)\cdot(p - p_3) =
	\det(p-p_1,p-p_2,p-p_3) = 0.\qedhere\]
\end{soln}

\begin{problem*}1.4.10.
	The natural orientation of $\RR^2$ makes it possible to associate
	a sign to the area $A$ of a parallelogram generated by
	two linearly independent vectors $u,v\in\RR^2$. To do this,
	let $\{e_1,e_2\}$ be the natural ordered basis of $\RR^2$,
	and write $u = u_1e_2 + u_2e_2$, $v = v_1e_1 + v_2e_2$.
	Observe the matrix relation
	\[\begin{pmatrix}
		u\cdot u & u\cdot v\\
		v\cdot u & v\cdot v
	\end{pmatrix} =
	\begin{pmatrix}
		u_1&u_2\\v_1&v_2
	\end{pmatrix}
	\begin{pmatrix}
		u_1&v_1\\u_2&v_2
	\end{pmatrix}\]
	and conclude that
	\[A^2 = \begin{vmatrix}
		u_1&u_2\\v_1&v_2
	\end{vmatrix}^2.\]
	Since the last determinant has the same sign as
	the basis $\{u,v\}$, we can say $A$ is positive or negative
	according to whether the orientation of $\{u,v\}$ is positive or negative.
	This is called the \vocab{oriented area} in $\RR^2$.
\end{problem*}

\begin{soln}
	The square of the area is
	\[\begin{vmatrix}
		u\cdot u & u\cdot v\\
		v\cdot u & v\cdot v
	\end{vmatrix} = \begin{vmatrix}
		u_1&u_2\\v_1&v_2
	\end{vmatrix}
	\begin{vmatrix}
	u_1&v_1\\u_2&v_2
	\end{vmatrix}
	= \begin{vmatrix}
		u_1&u_2\\v_1&v_2
	\end{vmatrix}^2,\]
	noting that a matrix has the same determinant as its transpose.
\end{soln}

\begin{problem*}1.4.11(a).
	Show that the volume $V$ of a parallelepiped generated by
	three linearly independent vectors $u,v,w\in\RR^3$ is given by
	$V = |(u\wedge v)\cdot w|$, and introduce an
	\vocab{oriented volume} in $\RR^3$.
\end{problem*}

\begin{soln}
	The expression $|(u\wedge v)\cdot w|$ is equal to
	$|u\wedge v|\cdot|w||\cos\theta|$, where $\theta$ is the angle formed by
	the vector $w$ and the normal to the $uv$-plane
	(since $u\wedge v$ is normal to the $uv$-plane).
	
	We know that the norm of $u\wedge v$ is the area of
	the parallelogram generated by $u$ and $v$, and $|w||\cos\theta|$
	is the length of the projection of $w$ on $u\wedge v$,
	i.e. the distance from $w$ to the $uv$-plane, which is
	the height of the parallelepiped between the two faces
	that are parallel to the $uv$-plane.
	Therefore, the volume of the parallelepiped is
	the product of the base area and the height, i.e. $|(u\wedge v)\cdot w|$.
\end{soln}






















\end{document}